\section{F02: certificate success can freeze precision refinement}\label{sec:f02}

\textbf{Classification: medium-priority accuracy-controller defect, not a false certificate.} The inspected exact interval proof is structurally sound for the supported expression class. The problem is that a successful proof of \emph{some} error bound can prevent the controller from increasing the arithmetic precision needed to prove the \emph{requested} error bound.

\source{AsymptoticInverse/Kernel/InverseCertificates.wl}{\code{certRound}, \code{certAttempt}, and the public \code{InverseCertificate} refinement loop.\cite{cert}}

\subsection{The proof already established by each successful attempt}

Let $I=[a,b]$ be a source interval, let $c\in(a,b)$ be a rational center, and suppose a structurally evaluated derivative enclosure $D$ has one sign and satisfies $|f'|\ge\mu>0$ on $I$. Let
\[
 f(c)-y\in R=[r_-,r_+],\qquad \varepsilon=\max(|r_-|,|r_+|).
\]
The mean-value radius is $r=\varepsilon/\mu$. If $[c-r,c+r]\subseteq I$, continuity and monotonicity establish a unique root there. When that containment fails, the implementation can instead prove opposite endpoint signs. Once existence is known, the mean-value identity also yields
\begin{equation}\label{eq:sharp-interval}
 x_*\in J=I\cap[c-r,c+r]\cap\left(c-\frac{R}{D}\right).
\end{equation}
The division is interval division with outward rational rounding. The implementation stores this sharper intersection as \code{RootEnclosure}. Decimal arithmetic chooses a seed; the successful proof uses rational endpoints and explicit elementary-function tail bounds.\cite{cert}

This is a useful trust boundary. Merely obtaining a high-precision \code{FindRoot} value would not supply the same proof object. Conversely, nothing here implies that native Wolfram interval facilities cannot be used to build similar proofs; the distinction is the package's specialized retained certificate and equation provenance.\cite{wl-interval,cert}

\subsection{The controller's two independent tasks are conflated}

The loop has to solve two different problems: establish a root enclosure at all, and make the center-error bound sufficiently small. Its source structure is essentially:
\begin{lstlisting}
result = certAttempt[..., order, ...];
If[certificateSucceeded && requestedAccuracyReached, return];
If[certificateSucceeded && centerWasAutomatic,
  interval = result["RootEnclosure"];
  center = Mean[interval];
  optionallyRefineTheAsymptoticSeed[],
  order = Min[2000, 2 order]
];
\end{lstlisting}
Thus an automatic-center attempt that succeeds but misses the accuracy goal narrows the source interval and changes the center, yet retains the same enclosure order. Arithmetic precision rises only in the other branch. Each attempt can continue to certify a true but inadequate bound at the same rounding resolution. Increasing \code{MaxRefinements} repeats the wrong type of work.\cite{cert}

The initial-order heuristic compounds this for relative-only requests: the decimal-digit estimate reads \code{TargetError} but not \code{RelativeError}. With no absolute tolerance, a very small relative tolerance does not affect the starting order. An adaptive controller could compensate for that omission, but the successful-attempt branch is precisely where compensation can be suppressed.\cite{cert}

\subsection{A restricted exact-rational experiment}

The supplied Python model specializes the relevant arithmetic to
\[
 f(x)=x^2,\quad y=2,\quad I_0=[1,2].
\]
It implements the inspected dyadic outward-rounding formula, the relevant interval products, the derivative/residual proof, and the intersection \eqref{eq:sharp-interval}. A 90-digit decimal approximation supplies an exact rational initial center. This seed is an experimental choice, not a claim about the native seed generated by the package. Every resulting root interval is independently checked by the exact inequalities
\[
 0<J_-\quad\text{and}\quad J_-^2\le2\le J_+^2.
\]

Set the initial enclosure order to $2$, so that the inspected bit rule gives $4(2+10)=48$ significant bits, and request absolute error $10^{-60}$. The baseline scheduling model performs 65 successful attempts without raising the order. Its final reported radius is exactly
\begin{equation}
 r_{64}=\frac{1}{199032864766430}
 \approx5.02429586778809\times10^{-15}.
\end{equation}
The requested accuracy is not reached. Under a simple alternative policy that doubles the order after every unsuccessful accuracy test, the same model reaches its goal in seven attempts, at order $128$ and 552 bits. The final reported radius is approximately $9.9643\times10^{-91}$.

These are exact-model results, not native package timings or native return values. The comparison establishes a concrete arithmetic/refinement mechanism and a repair direction; the delivered native characterization call separately tests the actual package:
\begin{lstlisting}
s = AsymptoticInverse[x^2, {x, 0}, {y, 2}];
InverseCertificate[s, 2,
  "Interval" -> {1, 2},
  "TargetError" -> 10^-60,
  "EnclosureOrder" -> 2,
  "MaxRefinements" -> 64,
  "RefineExpansion" -> False]
\end{lstlisting}
The expected source-derived symptom is an unmet accuracy request with unchanged enclosure order after successful automatic-center attempts. The exact native history remains to be measured.

\subsection{The stored intersection already proves a better error bound}

The controller also discards useful information when it reports only $r=\varepsilon/\mu$ as \code{CertifiedErrorBound}. Since $x_*\in J$, an additional rigorous upper bound is
\begin{equation}\label{eq:tight-error}
 B_J=\max\{|J_--c|,|J_+-c|\},\qquad
 B=\min\{\varepsilon/\mu,B_J\}.
\end{equation}
No new symbolic enclosure or function evaluation is required. A corresponding lower bound is the maximum of the existing residual-based lower bound and the distance from $c$ to $J$.

In the last baseline model attempt above, the root intersection proves
\[
 B_J=\frac{1}{5192296858534827628530496329220096}
 \approx1.92593\times10^{-34},
\]
while the reported radius remains approximately $5.0\times10^{-15}$. The code has already proved about nineteen more decimal orders of center accuracy than the controller uses. This particular sharper bound still does not reach $10^{-60}$, so both issues matter: tightening the reported bound is not a substitute for raising arithmetic precision.

\begin{proposition}
If a successful certificate establishes $x_*\in J$, replacing its center-error upper bound by \eqref{eq:tight-error} preserves certification. Raising the lower bound by the distance from $c$ to $J$ also preserves certification.
\end{proposition}
\begin{proof}
The continuous function $x\mapsto|x-c|$ attains its maximum over an interval at an endpoint and its minimum at $c$ if $c\in J$, otherwise at the nearest endpoint. Both interval-distance bounds therefore hold for every $x_*\in J$. Intersecting them with previously proved bounds cannot remove the actual error.
\end{proof}

The supplied \code{TightenCertificate} helper performs this post-processing on an already successful exact-rational certificate. It does not certify a new root or enclose an unspecified input-remainder family. It preserves the original equation and source-domain provenance.

\subsection{A controller that responds to the limiting error source}

A minimal safe repair is to raise the enclosure order after a successful attempt that misses its goal and fails a progress test. A better policy distinguishes the following limiting quantities: center displacement, width of the residual enclosure at a point, width of the derivative enclosure, and width of the verified root interval.

If the center is poor but arithmetic uncertainty is small, changing the center is productive. If the residual enclosure straddles zero largely because of rounding or elementary-function tail error, recomputing at higher enclosure order is productive. If the derivative is near zero or has severe interval dependency, source-interval narrowing can be productive. These actions are not mutually exclusive; the controller should not encode certificate success as a reason never to raise arithmetic precision.

A practical state transition can be specified as follows:
\begin{lstlisting}
certify at current interval, center, and enclosure precision;
if successful:
    tighten the error using the root intersection;
    retain the best proved certificate so far;
    if the requested mixed tolerance is proved: return;
    choose a new center inside the proved root enclosure;
if arithmetic uncertainty dominates or progress has stalled:
    increase enclosure order / significant bits;
if useful and allowed:
    refine the asymptotic seed;
stop with the best certificate and the reason for noncompletion;
\end{lstlisting}
The model's doubling-after-every-miss policy is deliberately simpler than this proposal. Its purpose is to demonstrate that the required accuracy is recoverable by changing the correct resource, not to prescribe optimal production scheduling.

For mixed absolute and relative accuracy, a proved lower bound $m\le|x_*|$ gives a sufficient absolute threshold
\[
 \tau=\max(\tau_{\rm abs},\tau_{\rm rel}m).
\]
The initial precision planner can use this threshold once $m>0$ is available. Before such a magnitude bound exists, it should retain the relative demand as an unresolved precision requirement rather than silently omit it. The existing safeguards concerning relative accuracy at a zero root should remain intact.\cite{cert}

The best-certificate store should compare proved bounds, not merely replace the previous result with the last association. The source currently records the latest successful attempt as \code{best}; there is no ordering test. A future enclosure-order change or seed replay need not monotonically improve the center-error bound, so retaining the best by an explicit criterion makes the failure result more useful without changing the proof system.

\subsection{Why the existing tests do not exercise this case}

The inspected acceptance tests sweep absolute and relative tolerances through $10^{-20}$ with a generously chosen enclosure order. Their explicit arithmetic-precision recovery test uses a \emph{fixed} center and an initial derivative-separation failure, which selects the controller branch that already doubles the order. They do not test a long sequence of successful but insufficient automatic-center certificates.\cite{certtests}

Add a deliberately low-order, automatic-center, tight-tolerance test; a relative-only high-accuracy test; and an assertion that the public error bound is no weaker than the bound computed from the retained root enclosure. Separate the statement ``a certificate was obtained'' from ``the requested accuracy was reached'' in both assertions and diagnostics. This is new coverage of a specific state-machine defect, not a repetition of the earlier general requests for better certificates or a nonempty native test gate.
